\section{GPT-3：Decoder-Only 与 In-Context Interface}

应用部分先用 GPT-3 展示 decoder-only 路线。模型堆叠 causal decoder blocks，以 next-token prediction 在大规模文本上预训练。训练目标简单，却把许多任务统一为“给定前缀，预测合理续写”。Prompt 中的 instruction、example 与 context 都成为同一 token stream。

\subsection{Autoregressive pretraining 与 prompting}

给定 token 序列 $x_1,\ldots,x_T$，语言模型最大化
\[
\sum_{t=1}^{T}\log p_\theta(x_t\mid x_{<t}).
\]
Causal mask 保证位置 $t$ 只看前缀。\term{in-context learning}（上下文学习）指模型在不更新参数的情况下，根据 prompt 内的指令或示例改变当前任务行为。它不是传统 gradient-based training，也不保证得到稳定规则；prompt format、example order 与 distribution shift 都会影响结果。

\lecturefigure{slide-16-gpt3.jpg}{GPT-3 用 decoder-only causal pretraining 与 prompt examples 形成通用生成接口。}{官方视频 18:30；Brown et al. 2020。}

\begin{knowledgebox}{读图：区分 pretraining、fine-tuning 与 in-context learning}
Pretraining 更新参数，让模型学习广泛 token distribution；fine-tuning 在任务数据上继续更新参数；in-context learning 只改变输入上下文，参数保持不变。图中的 few-shot 示例说明 task 被编码进 prompt，而不是证明模型真的执行了内部梯度下降。评估时要比较 zero-shot、few-shot、fine-tuned baseline，并控制 token budget 与 contamination。
\end{knowledgebox}

\teachervoice{讲者用 GPT-3 强调“把预训练网络当 generator”：给任务描述和少量 example，就能产生 completion。课堂提示是接口变化——许多任务从修改模型代码，转成设计上下文；这同时提高复用性，也把可靠性问题转移到 prompt、evaluation 与 guardrails。}

\begin{warningbox}{Next-token objective 不等于事实或意图对齐}
语言模型优化的是条件 token likelihood，不是“只说真实信息”或“完成用户真正目标”。规模和 prompting 可以提高任务表现，但 factuality、calibration、安全与 tool-use correctness 需要额外数据、训练和系统控制。
\end{warningbox}

\subsection{本章小结}

GPT-3 把 decoder-only Transformer 变成通用 prefix-to-continuation interface。它展示了预训练规模与 in-context behavior，却没有取消任务定义、评估和安全边界。

